\section{The energy obstacle map and its contact defect}
\label{sec:r12-energy-obstacles}

Keep exactly the strict original analytic Laurent box and background
extension \textup{(ED1)}.  Thus $S$ remains an arbitrary original
symmetric CLD middle-positive separately $d'$- and $d''$-closed
$(p,p)$-current, $q=n-p\ge1$, and $d=q+1$.  No factorization of
$S$ is imposed in this section.  Write
\[
 \mathscr C=\widehat{\mathcal E}^{1}_{\rm str}(V),\qquad
 \mathbb D_S=\mathscr C-\mathscr C,\qquad
 r(u)=e(u)^{1/d},\qquad E(u)=e(u)/d.
 \tag{OC1}
\]
The Banach norm $\|\cdot\|_{\rm ed}$, its literal Borel
realization, and the curvature functional
$J(u)(h)=-\int h\,d\widehat\mu_u$ are the proved R12
\textup{(ED2)--(ED11)}.  Their proof does not use the obstacle
map constructed below.  The R11 strong relative Perron theorem
\textup{(SE1)--(SE6)} is an \texttt{INTERNAL\_INPUT}.
All orders are relative Borel orders modulo local ordinary-open
outer $S$-capacity zero, denoted by $\preceq$.  All measures are
the original measures on the full original test sheaf.

\begin{definition}[An obstacle in the energy-difference space]
\label{def:r12-energy-obstacle}
For $f\in\mathbb D_S$ put
\[
 P(f)=\sup_{\rm rel}\{v\in\mathscr C:v\preceq f\},
 \qquad \mathfrak E(f)=E(P(f)).
 \tag{OC2}
\]
The supremum is the R11 relative least upper bound among Borel
classes.  It is not the ordinary upper semicontinuous
regularization of a pointwise supremum.
\end{definition}

\begin{theorem}[A global obstacle map and quantitative continuity]
\label{thm:r12-global-energy-obstacle}
The definition is valid for every $f\in\mathbb D_S$.  Its value
$P(f)$ is the greatest member of $\mathscr C$ below $f$, and
\[
 P(f)\preceq f,\qquad P(u)=u\quad(u\in\mathscr C),\qquad
 r(P(f))\le\|f\|_{\rm ed}.
 \tag{OC3}
\]
The map $P$ is order preserving, positively homogeneous and
superadditive in the relative order, hence order concave.
For $f,g\in\mathbb D_S$, set
$R=\max(r(P(f)),r(P(g)))$.  Then
\[
 \begin{split}
 |r(P(f))-r(P(g))|&\le\|f-g\|_{\rm ed},\\
 I(P(f),P(g))&\le R^q\|f-g\|_{\rm ed},\\
 \|J(P(f))-J(P(g))\|_*
       &\le\sqrt q\,R^{(2q-1)/2}\|f-g\|_{\rm ed}^{1/2}.
 \end{split}
 \tag{OC4}
\]
At $R=0$ the second and third left sides are zero.  In
particular, norm convergence of obstacles gives strong
convergence of their envelopes and dual-norm convergence of
their original curvature functionals.  Norm continuity of
$P$ with values in $\mathbb D_S$ is not claimed.
\end{theorem}

\noindent\textbf{Status: \texttt{NEW\_PROVED}} for the
global energy-difference obstacle domain and its quantitative
estimates, in every degree and for arbitrary original $S$.

\begin{proof}
Write $f=a-b$ with $a,b\in\mathscr C$.  Since $b\preceq0$,
one has $a\preceq f$.  The family in \textup{(OC2)} is therefore
nonempty.  R11 \textup{(SE2)} gives its relative envelope in
$\mathscr C$ and keeps it below the Borel upper constraint $f$.
Thus the envelope itself belongs to the family and is its
greatest member.  In particular $P(u)=u$ on the cone.
Energy monotonicity gives $r(P(f))\le r(a)$; infimizing over
decompositions proves the last assertion of \textup{(OC3)}.

Increasing the obstacle increases its competitor family.
Positive rescaling maps competitor families bijectively onto
each other, and $P(0)=0$ handles scale zero.  Also
$P(f)+P(g)\in\mathscr C$ lies below $f+g$, proving
superadditivity and hence concavity.

For the estimates write $f-g=a-b$, $a,b\in\mathscr C$, and
put $u=P(f)$, $v=P(g)$.  Since
\[
 u+b\preceq f+b=g+a\preceq g,\qquad
 v+a\preceq g+a=f+b\preceq f,
\]
both are legitimate cone competitors.  Consequently
\[
                        a\preceq u-v\preceq-b.
 \tag{OC5}
\]
In particular $u\succeq v+a$ and $v\succeq u+b$.
Energy monotonicity and Minkowski give
$r(u)\le r(v)+r(a)$ and $r(v)\le r(u)+r(b)$.
Infimizing proves the first estimate in \textup{(OC4)}.

The literal mixed integrals and \textup{(OC5)} give the sharper
signed estimate, without inserting an absolute-variation norm,
\[
 \begin{split}
 I(u,v)&=J(u)(u-v)-J(v)(u-v)\\
       &\le J(u)(a)+J(v)(b)\\
       &\le r(a)r(u)^q+r(b)r(v)^q
        \le R^q(r(a)+r(b)).
 \end{split}
 \tag{OC6}
\]
Taking the infimum gives the second estimate.  Combine it with
the proved integrated Cauchy--Schwarz bound
$\|J(u)-J(v)\|_*\le\sqrt q R^{(q-1)/2}I(u,v)^{1/2}$
to obtain the third.  The bound in \textup{(OC3)} controls
$R$ on a norm-convergent obstacle row.  All subsequent
strong current and weighted limits are therefore the full
original limits of R11, not only selected moment pairings.
\end{proof}

\begin{proposition}[A finite convex obstacle energy]
\label{prop:r12-convex-obstacle-energy}
The function $\mathfrak E:\mathbb D_S\to[0,\infty)$ is
finite, convex, order decreasing and positively homogeneous
of degree $d$.  It is locally Lipschitz, with the precise bound
\[
       |\mathfrak E(f)-\mathfrak E(g)|
           \le\max(r(P(f)),r(P(g)))^q\|f-g\|_{\rm ed}.
 \tag{OC7}
\]
Its restriction to $\mathscr C$ is the original first energy
$e/d$.  Its convexity is a statement on all of $\mathbb D_S$;
it does not declare the polynomial extension $B(f^{[d]})/d$
convex outside the potential cone.
\end{proposition}

\noindent\textbf{Status: \texttt{NEW\_PROVED}}.

\begin{proof}
Finiteness, order decrease and homogeneity follow from
\textup{(OC3)} and the corresponding cone energy properties.
For $0\le s\le1$, the order concavity of $P$ gives
\[
 P((1-s)f+sg)\succeq(1-s)P(f)+sP(g).
\]
Apply energy monotonicity, followed by convexity of $E$ on
the cone, to obtain convexity of $\mathfrak E$.  Finally,
$|x^d-y^d|/d\le\max(x,y)^{d-1}|x-y|$ for $x,y\ge0$.
Use the first line of \textup{(OC4)} and $d-1=q$ to obtain
\textup{(OC7)}.  The norm bound \textup{(OC3)} makes this a
local Lipschitz estimate on the entire Banach space.
\end{proof}

\subsection{Contact defects and exact subgradients}

Call a continuous linear functional $\ell\in\mathbb D_S^*$
\emph{order reversing} if $h\succeq0$ implies $\ell(h)\le0$.
This is order on Borel scalar differences, not an identification
of positivity cones for differential forms.  Every literal
positive measure datum $\ell_\mu(h)=-\int h\,d\mu$ with finite
integrals and a continuous extension to $\mathbb D_S$ has this
property.  Write
\[
 \partial\mathfrak E(f)=\{\ell\in\mathbb D_S^*:
 \mathfrak E(g)\ge\mathfrak E(f)+\ell(g-f)
                      \text{ for every }g\in\mathbb D_S\}
 \tag{OC8}
\]
for the Banach subdifferential, and retain the cone polar
$\mathscr C^\circ$ of \textup{(NI5)}.

\begin{definition}[Relative contact defect]
\label{def:r12-relative-contact-defect}
For $f\in\mathbb D_S$, with $u=P(f)$, put
\[
 \mathfrak Q(f)=\int_V(f-u)\,d\widehat\mu_u
                         =J(u)(u-f)\ge0.
 \tag{OC9}
\]
The integral is literal and finite: $f-u$ is a finite sum of
cone differences, and all its absolute mixed first-energy
integrals are finite.  Its nonnegativity uses $u\preceq f$
and the fact that the original completed measure ignores the
relative null set.  Vanishing means the relative Borel
contact identity
$(f-u)\widehat\mu_u=0$; no ordinary support regularity is
part of this definition.
\end{definition}

\begin{theorem}[The obstacle subdifferential and contact criterion]
\label{thm:r12-obstacle-subdifferential}
Let $f\in\mathbb D_S$ and $u=P(f)$.  A functional
$\ell\in\mathbb D_S^*$ belongs to $\partial\mathfrak E(f)$
if and only if all three conditions hold:
\[
 \begin{gathered}
 \ell\text{ is order reversing},\\
 \ell-J(u)\in\mathscr C^\circ,
                   \qquad (\ell-J(u))(u)=0,\\
                         \ell(f-u)=0.
 \end{gathered}
 \tag{OC10}
\]
Every such functional has $\|\ell\|_*=r(u)^q$.
In particular,
\[
       \boxed{\quad J(P(f))\in\partial\mathfrak E(f)
                        \quad\Longleftrightarrow\quad
                         \mathfrak Q(f)=0.\quad}
 \tag{OC11}
\]
For $u\in\mathscr C$, one always has
$J(u)\in\partial\mathfrak E(u)$.

If $\ell$ is additionally in the strong-continuous dual
$\mathbb D_{S,\mathrm s}^*$, the middle line of
\textup{(OC10)} is equivalent to $u=\Pi(\ell)$ by the
proved R12 normal-decomposition theorem.  Strong continuity
is not being asserted for every Banach subgradient.
\end{theorem}

\noindent\textbf{Status: \texttt{NEW\_PROVED}} for the
precise original arbitrary-background obstacle energy.  This is
a contact/variational characterization; it does not assert
that the contact defect vanishes.

\begin{proof}
Suppose first that $\ell\in\partial\mathfrak E(f)$.
For $h\succeq0$ and $t>0$, order decrease gives
$\mathfrak E(f+th)\le\mathfrak E(f)$.  The subgradient
inequality then gives $t\ell(h)\le0$, proving order reversal.
Since $P(u)=u$ and $E(u)=\mathfrak E(f)$, the competitor
$g=u$ in \textup{(OC8)} gives $\ell(u-f)\le0$.
But $u-f\preceq0$ and order reversal give the opposite
inequality.  Hence $\ell(f-u)=0$.

Restricting \textup{(OC8)} to $v\in\mathscr C$ now gives
\[
                        E(v)\ge E(u)+\ell(v-u).
 \tag{OC12}
\]
Along the cone segment $u+t(v-u)$ the polynomial energy
has derivative $J(u)(v-u)$ at zero.  Applying
\textup{(OC12)}, dividing by $t>0$ and letting $t\downarrow0$
gives $(\ell-J(u))(v-u)\le0$.  The competitors $0$ and
$2u$ yield $(\ell-J(u))(u)=0$, and then the same inequality
is the polar condition.  This proves necessity.

Conversely assume \textup{(OC10)}.  Cone convexity and its
first derivative give
\[
 E(v)\ge E(u)+J(u)(v-u)\ge E(u)+\ell(v-u)
                                     \quad(v\in\mathscr C).
\]
For arbitrary $g\in\mathbb D_S$ use $v=P(g)\preceq g$.
Order reversal gives $\ell(P(g)-u)\ge\ell(g-u)$, so
\[
 \mathfrak E(g)\ge E(u)+\ell(g-u)
                 =\mathfrak E(f)+\ell(g-f).
\]
The equality uses $\ell(f-u)=0$.  This proves sufficiency.

Order reversal makes $\ell(v)\ge0$ on the nonpositive cone.
The polar condition gives
$0\le\ell(v)\le J(u)(v)\le r(v)r(u)^q$.
The exact dual norm formula \textup{(ED10)} proves
$\|\ell\|_*\le r(u)^q$.  If $u\ne0$, evaluation on
$u/r(u)$ gives equality because $\ell(u)=e(u)$.
For $u=0$, the upper bound gives the same conclusion.
Finally $J(u)$ is order reversing, its normal residual is
zero, and its last condition in \textup{(OC10)} is exactly
$-\mathfrak Q(f)=0$.  This proves \textup{(OC11)} and
the assertion at cone points.  The final equivalence is
precisely the already proved normal-decomposition theorem
with its stated continuity gate.
\end{proof}

\begin{corollary}[Zero-mass slack does not move a cone envelope]
\label{cor:r12-zero-mass-obstacle-slack}
If $u\in\mathscr C$, $g\in\mathbb D_S$, $g\succeq0$, and
$\int g\,d\widehat\mu_u=0$, then
$P(u+g)=u$ and $\mathfrak Q(u+g)=0$.
\end{corollary}

\noindent\textbf{Status: \texttt{DIRECT\_COROLLARY}} of
the frozen ordered energy separation and the new obstacle definition.
It is recorded as a valid nontrivial contact mechanism, not as
unrestricted orthogonality.

\begin{proof}
Put $v=P(u+g)$.  Then $u\preceq v\preceq u+g$.
Thus $\int(v-u)\,d\widehat\mu_u=0$ and convexity gives
$E(v)\ge E(u)+J(u)(v-u)=E(u)$.  Energy monotonicity gives
the reverse inequality.  The frozen ordered discrepancy bound
$I(u,v)\le e(u)-e(v)$ makes $u=v$ as relative classes.
The contact defect is then the assumed zero integral.
\end{proof}

\subsection{Directional loss, source equations and failure of orthogonality}

For $u\in\mathscr C$ and $h\in\mathbb D_S$, define the
one-sided defect
\[
 \mathfrak d_u(h)=
 \lim_{t\downarrow0}
       \frac{\mathfrak E(u+th)-E(u)}{t}-J(u)(h).
 \tag{OC13}
\]
Convexity and local Lipschitz continuity make the limit finite.
Since $J(u)$ is a subgradient, the defect is nonnegative.
It is positively homogeneous and subadditive, being a convex
directional derivative minus a linear functional.

\begin{theorem}[A first-order contact gate and its exact loss]
\label{thm:r12-first-order-contact-gate}
For $u,h$ as above, put $v_t=P(u+th)$ and
$G_t=u+th-v_t\succeq0$ for $t>0$.  Then
\[
 \begin{split}
 &\mathfrak E(u+th)-E(u)-tJ(u)(h)\\
 &\quad=
 \underbrace{E(v_t)-E(u)-J(u)(v_t-u)}_{\ge0}
                       +\int_VG_t\,d\widehat\mu_u\ge0,
 \end{split}
 \tag{OC14}
\]
and
\[
 \mathfrak E(u+th)-E(u)-tJ(u)(h)
       \le t(J(v_t)-J(u))(h)+\mathfrak Q(u+th).
 \tag{OC15}
\]
Consequently
\[
 \mathfrak d_u(h)
       \le\liminf_{t\downarrow0}\frac{\mathfrak Q(u+th)}{t}.
 \tag{OC16}
\]

Let $\mu$ be a positive Radon datum whose literal integration
functional $\ell_\mu\in\mathbb D_{S,\mathrm s}^*$ is defined
on every energy difference, and let $u=\Pi(\ell_\mu)$.
If, for every original compact smooth scalar $\phi$ and both
signs, one has
\[
                  \mathfrak Q(u\pm t\phi)=o(t)
                                     \quad(t\downarrow0),
 \tag{OC17}
\]
then $\widehat\mu_u=\mu$ on every original compact smooth
scalar test, after every actual moment refinement.
Condition \textup{(OC17)} is a sufficient first-order contact
criterion, not a condition proved for general obstacles or data.
\end{theorem}

\noindent\textbf{Status: \texttt{NEW\_PROVED}} for the
defect identities and \texttt{CONDITIONAL} for the source
conclusion under the explicit geometric contact criterion.
This condition is stated using the independently constructed
envelope and its own curvature, not by assuming the desired
source equality.

\begin{proof}
All integrals are finite by the mixed first-energy inequality.
Subtract $tJ(u)(h)$ from the convex tangent remainder at $u$.
Since $v_t-u-th=-G_t$ and $J(u)=-\int(\cdot)\,d\widehat\mu_u$,
this gives exactly \textup{(OC14)}.  The tangent inequality
at $v_t$ instead gives
$E(v_t)-E(u)\le J(v_t)(v_t-u)$.
Substitute $v_t-u=th-G_t$ to obtain \textup{(OC15)}.
The third estimate of \textup{(OC4)} gives
$J(v_t)\to J(u)$ in the actual dual norm.  Divide
\textup{(OC15)} by $t$ and let $t\downarrow0$; the error
$(J(v_t)-J(u))(h)$ tends to zero.  This proves
\textup{(OC16)}, including the possibility that its right
side is infinite.

For the datum assertion, order reversal and the normal
decomposition give $\ell_\mu\in\partial\mathfrak E(u)$
by \textup{(OC10)}.  Thus
\[
 \mathfrak d_u(h)\ge(\ell_\mu-J(u))(h)
                                    \quad(h\in\mathbb D_S).
 \tag{OC18}
\]
Every original compact smooth $\phi$ belongs to $\mathbb D_S$
by the exact all-degree finite-mass difference construction
R11 \textup{(HC14)--(HC15)}.  Assumption \textup{(OC17)},
\textup{(OC16)} and nonnegativity make both
$\mathfrak d_u(\phi)$ and $\mathfrak d_u(-\phi)$ zero.
Apply \textup{(OC18)} in both signs to get
$\ell_\mu(\phi)=J(u)(\phi)$.  The literal integral
definitions give the original source identity on every
such test.  The choice of $\phi$ was unrestricted in the
original sheaf, so actual refinements are included.
\end{proof}

\begin{corollary}[A nonzero variational residual forces an obstacle kink]
\label{cor:r12-normal-residual-obstacle-kink}
For the positive datum and maximizer of the preceding theorem,
put $n=\ell_\mu-J(u)$.  For every $h\in\mathbb D_S$ the
right and left directional slopes of $\mathfrak E$ at $u$
have a gap at least $|n(h)|$.  If $n(h)>0$, then
\[
 \mathfrak d_u(h)\ge n(h)>0,
 \qquad
 \liminf_{t\downarrow0}\frac{\mathfrak Q(u+th)}{t}\ge n(h).
 \tag{OC19}
\]
If the source equation fails on an original compact smooth
test, the direction $h$ can be chosen to be such a test,
with its sign changed if necessary.  In that case universal
relative obstacle orthogonality and two-sided differentiability
of this obstacle energy at the maximizer are both false.
\end{corollary}

\noindent\textbf{Status: \texttt{NEW\_PROVED}} as a
quantitative contact interpretation of the normal residual.

\begin{proof}
Both $J(u)$ and $\ell_\mu$ are subgradients at $u$.
The right slope in direction $h$ is at least the larger of
$J(u)(h)$ and $\ell_\mu(h)$, whereas the left slope is at
most their smaller value.  Their difference is at least
$|n(h)|$.  Equation \textup{(OC18)} gives the first bound
in \textup{(OC19)}, and \textup{(OC16)} gives the second.
Distinct original Radon measures are distinguished by an
original compact smooth scalar test; the exact test-difference
construction places that direction in $\mathbb D_S$.
No ordinary upper envelope or stronger delta-test equality
is used in this separation.
\end{proof}

\paragraph{Which proposed classical step fails.}
The map $P$ exists on every energy difference, its energy is
finite and convex, and the original curvature is strongly
continuous.  These facts alone do not make the curvature a
subgradient at the \emph{obstacle} $f$: by \textup{(OC11)},
the missing statement is precisely its nonnegative contact
defect.  The new all-degree compact-source examples force
that defect to persist at first order along an original
smooth direction through a genuine variational maximizer.
Thus the unrestricted contact step is \texttt{FALSE} in
this classical-approximation strong class.  The qualified
criterion \textup{(OC17)} remains useful on a separately
proved obstacle/data subclass; no such universal subclass
identification is claimed.
